% unidades/26-condicional-1.tex
\chapter{🔀 条件① — Condicional: と・ば}
\unidadheader{26}{条件①}{Condicional と・ば}{Expresar consecuencias naturales e hipótesis}

\begin{tcolorbox}[vocabbox, title={📌 Vocabulario Nuevo}]
\begin{tabularx}{\linewidth}{llX}
\toprule
\textbf{Japonés} & \textbf{Español} & \textbf{Tipo} \\
\midrule
\vocabitem{押す}{おす}{presionar / empujar}{v. gr.1}
\vocabitem{引く}{ひく}{jalar / tirar}{v. gr.1}
\vocabitem{開く}{あく}{abrirse (intrans.)}{v. gr.1}
\vocabitem{閉まる}{しまる}{cerrarse (intrans.)}{v. gr.1}
\vocabitem{動く}{うごく}{moverse}{v. gr.1}
\vocabitem{止まる}{とまる}{detenerse}{v. gr.1}
\vocabitem{曲がる}{まがる}{doblar / girar}{v. gr.1}
\vocabitem{まっすぐ}{—}{recto / directo}{adv.}
\vocabitem{角}{かど}{esquina}{sust.}
\vocabitem{〜と〜}{—}{si ~ entonces ~ (natural)}{conj.}
\bottomrule
\end{tabularx}
\end{tcolorbox}

\subsection*{🗣️ Frases Clave}

\frase{このボタンを\ruby{押}{お}すと、\ruby{開}{あ}きますよ。}
  {Si presionas este botón, se abrirá.}
\frase{まっすぐ\ruby{行}{い}くと、\ruby{駅}{えき}がありますよ。}
  {Si vas recto, encontrarás la estación.}
\frase{お\ruby{金}{かね}があれば、\ruby{旅行}{りょこう}したいです。}
  {Si tuviera dinero, quisiera viajar.}

\subsection*{⚙️ Gramática}

\patron{Consecuencia natural: と}
  {普通形 + と + 結果 (consecuencia automática / instrucciones)}
  {Usado para: instrucciones, hechos naturales, rutas.
  NO para: deseos, peticiones, invitaciones.
  \ej{この\ruby{薬}{くすり}を\ruby{飲}{の}むと、\ruby{眠}{ねむ}くなります。}
    {Si tomas esta medicina, te dará sueño.}
  \ej{\ruby{右}{みぎ}に\ruby{曲}{ま}がると、\ruby{銀行}{ぎんこう}があります。}
    {Si doblas a la derecha, hay un banco.}}

\patron{Hipótesis: ば}
  {V-えば / adj-ければ / N-であれば}
  {Más formal y literario que と. Expresa condición hipotética.
  \ej{\ruby{春}{はる}になれば、\ruby{桜}{さくら}が\ruby{咲}{さ}きます。}
    {Cuando llegue la primavera, florecerán los cerezos.}
  \ej{\ruby{時間}{じかん}があれば、\ruby{手伝}{てつだ}います。}
    {Si tengo tiempo, te ayudaré.}
  \ej{\ruby{高}{たか}くなければ、\ruby{買}{か}います。}
    {Si no es caro, lo compraré.}}

\begin{tcolorbox}[ejerciciobox, title={📝 Ejercicios Rápidos}]
\begin{enumerate}
  \item Usa と para dar instrucciones de cómo llegar a algún lugar.
  \item ¿と o ば? Esta oración: 春\_\_なる\_\_、暖かくなります。
  \item Traduce: \ruby{右}{みぎ}に\ruby{曲}{ま}がれば、\ruby{学校}{がっこう}が\ruby{見}{み}えます。
\end{enumerate}
\vspace{0.4em}
\textbf{Respuestas:}
1) (personal)\quad
2) と (\ruby{春}{はる}になると、\ruby{暖}{あたた}かくなります。)\quad
3) Si doblas a la derecha, verás la escuela.
\end{tcolorbox}

\begin{tcolorbox}[kanjibox, title={🀄 Kanjis de la Unidad}]
\begin{tabularx}{\linewidth}{cllXX}
\toprule
\textbf{Kanji} & \textbf{On} & \textbf{Kun} & \textbf{Significado} & \textbf{Vocab} \\
\midrule
\kanjirow{押}{おう}{お}{presionar}{\ruby{押}{お}す / \ruby{押入}{おしい}れ}
\kanjirow{引}{いん}{ひ}{jalar}{\ruby{引}{ひ}く / \ruby{引越}{ひっこ}し}
\kanjirow{動}{どう}{うご}{mover(se)}{\ruby{動}{うご}く / \ruby{運動}{うんどう}}
\kanjirow{止}{し}{と・や}{detener(se)}{\ruby{止}{と}まる / \ruby{中止}{ちゅうし}}
\bottomrule
\end{tabularx}
\end{tcolorbox}
